\section{Conclusion}
\label{sec:conclusion}

We studied how transformer layers evolve representations by decomposing
learned transformations into parallel and perpendicular components. Across
pretrained models, learned attention and MLP outputs contain substantial
direction-preserving parallel components, showing that rescaling-like behavior
is not limited to the additive identity path. Their robustness is
space-dependent: manipulating the attention aggregation aligned with the
projected self-value remains closest to baseline, whereas residual-space
scaling is more fragile and perpendicular scaling is more disruptive.

The attention-side factorization separates diagonal self-routing, value
self-alignment, and cross-token aggregation, explaining why value-space and
residual-space edits can produce different behavior. The same geometry also
supports downstream analysis: stronger compressed models have lower
perpendicular transformation error, and retained pretraining runs indicate
that suppressing self-value-parallel components can shift optimization and
improve average downstream performance. Overall, parallel/perpendicular
geometry provides a compact way to analyze not only the size of a layer
transformation, but also whether it preserves or changes representation
direction.

\section{Limitations}
\label{sec:limitations}

Our analysis and experiments focus on the models, language tasks, and
intervention settings studied here. Although the same geometric picture is consistent
across the evaluated models, its behavior under other architectures, training
regimes, or multimodal task distributions remains to be established. We also
treat the proposed edits primarily as analysis tools; turning them into
practical training or inference methods requires further validation.
